% W5i82_NewFrontiers.tex
% DFD 20-Agent Research Programme --- Iteration 82
% New Frontiers, Paper Proposals, and Research Directions

\documentclass[11pt,a4paper]{article}
\usepackage[utf8]{inputenc}
\usepackage{amsmath,amssymb}
\usepackage{booktabs}
\usepackage{geometry}
\usepackage{xcolor}
\usepackage{enumitem}
\geometry{margin=2.5cm}

\title{W5i82: New Frontiers and Paper Proposals\\[6pt]
\large Textbook Launch $\cdot$ DFD Classification Theory $\cdot$ 2177 Total Proposals}
\author{DFD 20-Agent Research Programme}
\date{2026-03-24}

\begin{document}
\maketitle

\section{Paper Proposals (i82 Batch: 100 new)}

\subsection{Tier 1: High Priority (3 proposals)}

\begin{enumerate}[nosep]
\item \textbf{``Gravitational Wave Speed in DFD: $c_T = c$ Exactly''} (6 pp, CQG/PRD Rapid). A short paper proving $c_T = c$ in DFD, with implications for GW170817 and next-generation detectors. Can be submitted quickly. High impact because many modified gravity papers are eliminated by $c_T$ constraints.

\item \textbf{``Density Field Dynamics: From First Principles to Cosmological Predictions''} (Textbook, 400--500 pp, CUP/Springer). The definitive reference for DFD. Self-contained graduate-level text.

\item \textbf{``DFD Predictions for SPHEREx BAO and Galaxy Clustering''} (15--18 pp, MNRAS). SPHEREx launches Q3 2026. DFD predictions for its BAO and galaxy power spectrum measurements. Timely paper.
\end{enumerate}

\subsection{Tier 2: Medium Priority (12 proposals)}

\begin{enumerate}[nosep]
\setcounter{enumi}{3}
\item DFD predictions for DES Y6 cosmic shear
\item DFD in the EFTofDE language: mapping and comparison
\item DFD and the Alcock--Paczynski effect: geometric test
\item DFD predictions for Planck PR4 lensing
\item DFD void lensing predictions
\item DFD and photometric redshift calibration: does $\psi$ bias photo-$z$?
\item DFD predictions for SKA HI intensity mapping
\item DFD and the marked correlation function
\item DFD predictions for CMB $B$-mode delensing
\item DFD and the cosmic dipole
\item DFD predictions for FRB cosmology
\item DFD and the galaxy--void cross-correlation
\end{enumerate}

\subsection{Tier 3: Exploratory (85 proposals)}

Proposals 16--100: DFD predictions for various next-generation experiments and theoretical extensions including: strong lensing time delay cosmography, galaxy rotation curves (MOND comparison), ultra-high-energy cosmic ray propagation, dark matter substructure statistics, primordial magnetic field generation, cosmic string bounds, phase transitions in the early universe, topological defects in the $\psi$-field, higher-order statistics (Minkowski functionals), machine learning emulators for DFD, $\psi$-field and inflation (reheating), DFD and the swampland conjectures, DFD and de Sitter entropy, and many more.

\section{DFD Classification: Where Does DFD Sit?}

A key theoretical advance of i82 is the precise classification of DFD within the landscape of gravitational theories.

\begin{center}
\begin{tabular}{lccccc}
\toprule
\textbf{Theory} & \textbf{Extra DOF} & $c_T = c$ & \textbf{Screening} & \textbf{$\gamma_{\rm PPN} = 1$} & \textbf{Type} \\
\midrule
GR & 0 & Yes & N/A & Yes & Metric \\
\textbf{DFD} & \textbf{0} & \textbf{Yes} & \textbf{None needed} & \textbf{Yes} & \textbf{Constrained} \\
$f(R)$ & 1 scalar & No (generic) & Chameleon & No & Scalar-tensor \\
BD & 1 scalar & No (generic) & None & No & Scalar-tensor \\
DGP & 1 scalar & Yes & Vainshtein & No & Braneworld \\
Horndeski & 1 scalar & Tunable & Various & Tunable & Scalar-tensor \\
Massive gravity & 5 & No & Vainshtein & No & Spin-2 \\
\bottomrule
\end{tabular}
\end{center}

DFD occupies a unique position: it modifies gravity without adding propagating degrees of freedom, without requiring screening, and while maintaining $c_T = c$ and $\gamma_{\rm PPN} = 1$ exactly. The only other theory sharing all these properties is GR itself.

The difference between DFD and GR is the algebraic $\psi$-field constraint, which modifies the Poisson equation on cosmological scales while leaving local physics unchanged. This is the precise sense in which DFD is ``minimally modified gravity.''

\section{Textbook: Detailed Chapter Outlines}

\subsection{Chapter 1: The Density Field Principle (30 pp)}
\begin{itemize}[nosep]
\item 1.1 Motivation: why modify gravity?
\item 1.2 The density field as a fundamental object
\item 1.3 Historical context: Mach's principle, Dirac's large numbers
\item 1.4 The DFD postulates
\item 1.5 Overview of the theory's predictions
\item 1.6 How to read this book
\item Exercises (10 problems)
\end{itemize}

\subsection{Chapter 4: The Fine-Structure Constant (50 pp)}
\begin{itemize}[nosep]
\item 4.1 The $\alpha$ problem in physics
\item 4.2 Heat kernel methods on $\mathbb{C}P^2$
\item 4.3 The $a_4$ coefficient: complete derivation
\item 4.4 Dimensional transmutation and the $\alpha$ formula
\item 4.5 Numerical evaluation: $\alpha_{\rm DFD}$ vs.\ $\alpha_{\rm exp}$
\item 4.6 The residual: 0.55 ppm and what it means
\item 4.7 Common objections and responses
\item Exercises (20 problems)
\end{itemize}

\subsection{Chapter 7: Observational Predictions I (50 pp)}
\begin{itemize}[nosep]
\item 7.1 CMB temperature anisotropies
\item 7.2 CMB polarisation and $B$-modes
\item 7.3 Matter power spectrum
\item 7.4 BAO scale
\item 7.5 Growth rate $f\sigma_8$
\item 7.6 Weak gravitational lensing
\item 7.7 $E_G$ statistic
\item 7.8 Cluster counts and the halo mass function
\item 7.9 Comparison tables: DFD vs.\ data
\item Exercises (15 problems)
\end{itemize}

\section{Updated Paper Portfolio}

\begin{center}
\begin{tabular}{lccl}
\toprule
\textbf{Paper} & \textbf{Status} & \textbf{i82 Change} & \textbf{Next Step} \\
\midrule
$C_\ell$ paper & Submitted & --- & Awaiting referee \\
$\alpha$ PRL & 100\% & --- & Submit i83 \\
v3.3 unified & 100\% & --- & Submit i83 \\
DESI standalone & 100\% & --- & Phase 2 batch \\
Review paper & 100\% & --- & Phase 2 batch \\
No-screening & 100\% & --- & Phase 2 batch \\
Software/JOSS & 100\% & --- & Phase 2 batch \\
Fermion masses & 100\% & --- & Phase 2 batch \\
$E_G$ paper & \textbf{100\%} & \textbf{+5\%} & Phase 2 batch \\
Multi-messenger & 92\% & +7\% & i83 target: 100\% \\
CMB-S4 forecast & 70\% & --- & i84 push \\
N-body I & 22\% & +7\% & i84--i88 \\
BD letter & 45\% & +40\% & i83 target: 80\% \\
$c_T$ rapid comm. & \textbf{New: 30\%} & New & i83 target: 70\% \\
21\,cm paper & 5\% & --- & i85+ \\
SPHEREx paper & \textbf{New: 5\%} & New & i84+ \\
Textbook & \textbf{Outline} & New & Ch.\ 1 draft i83 \\
\bottomrule
\end{tabular}
\end{center}

\textbf{Total active papers}: 17 (9 at 100\%, 1 submitted, 7 in progress). Plus textbook.

\section{i83 Targets}

\begin{enumerate}[nosep]
\item N-body: 55\% $\to$ 65\%. Phase 2 Run A to 75\%.
\item Multi-messenger: 92\% $\to$ 100\% (10th paper at 100\%).
\item BD letter: 45\% $\to$ 80\%.
\item $c_T$ rapid communication: 30\% $\to$ 70\%.
\item Comprehensive literature update: What's new in 2026?
\item N-body paper I: 22\% $\to$ 35\%.
\item Textbook Chapter 1: first draft.
\end{enumerate}

\end{document}
